\documentclass[11pt]{article}

\usepackage[left=1in, right=2.1in, top=1in, bottom=1in, marginparwidth=1.65in]{geometry}
\usepackage[T1]{fontenc}
\usepackage{lmodern}
\usepackage{microtype}
\usepackage{amsmath,amssymb}
\usepackage{enumitem}
\usepackage{booktabs}
\usepackage{xcolor}
\usepackage[textsize=scriptsize]{todonotes}
\usepackage[hidelinks]{hyperref}

% Passages drafted with AI assistance and not yet checked against their sources.
% List them all with: grep -n unverified whitepaper.tex
\newcommand{\unverified}[1]{%
  \todo[color=orange!25, bordercolor=orange!70!black, linecolor=orange!70!black]{%
    \textbf{Unverified (AI-drafted).} #1}}

\title{From Relation to Node:\\[0.3em]
\large Minimal Conditions for the Emergence of Grounded Symbolic Graphs,\\
with a Machine-Verified Formalization in Lean~4}
% TODO: author name
\author{}
\date{Draft, \today}

\begin{document}

\maketitle

\begin{center}
\fbox{\parbox{0.92\linewidth}{\small
\textbf{Status.} This document is a work in progress. It is incomplete, and all of its
parts, including the axioms, definitions, results and their interpretation, are subject
to revision and refinement. Passages marked in the margin were drafted with AI
assistance and have not yet been checked against their sources.}}
\end{center}

\vspace{1em}
\begin{flushright}
\emph{Mysteries should not be multiplied beyond necessity.}\\[0.2em]
--- Eric Gans, \emph{The Origin of Language} (1981)
\end{flushright}

% TODO: abstract, once the later sections are drafted.

\section{Motivation and Background}

\subsection{Theories of language origin}

Theories of language origin attempt to explain how a species without language came to
have it. The field is broad. Among its main families are the following.
\begin{itemize}[nosep]
  \item Gestural accounts, on which language grew out of manual communication and
    pantomime \cite{Corballis2002,Arbib2012}.
  \item Social accounts, on which language took over the role of grooming in maintaining
    group bonds \cite{Dunbar1996}.
  \item Intention-based accounts, which ground meaning in the expression and recognition
    of communicative intentions \cite{Grice1957,Tomasello2008}.
  \item Game-theoretic accounts, in which signals acquire meaning as stable conventions
    among interacting agents \cite{Lewis1969,Skyrms2010}.
  \item Computational accounts, which locate the decisive step in the appearance of a
    single combinatorial operation \cite{BerwickChomsky2016}.
\end{itemize}
Two problems recur across these families, and they help explain why the question of
language origin has often been treated with suspicion.

\paragraph{Presupposition.}
Many accounts take as given, in some form, capacities closely related to the ones whose
origin is at issue. This is not always an oversight; often it is a deliberate choice of
what is to be explained. Three examples illustrate the point.

\begin{itemize}
  \item \emph{Game-theoretic accounts.} Lewis--Skyrms signaling games show how signals
    can acquire meaning without design or prior agreement. Skyrms extends the framework
    so that agents begin with no signals at all and invent new ones through a
    reinforcement process \cite{Skyrms2010,AlexanderSkyrmsZabell2012}. What these models
    take as given is the interaction itself: a sender whose behaviour depends on a state
    of the world, a receiver whose action depends on the sender's behaviour, and payoffs
    that reward a match. The models explain which behaviour comes to mean which state.
    That one participant's behaviour is the kind of thing the other's action responds to
    \emph{as} a signal is the frame within which they work.
  \item \emph{Intention-based accounts.} These analyze meaning through the expression
    and recognition of communicative intentions \cite{Grice1957}. As analyses of
    meaning, they are not committed to any story of origin. As a basis for one, they
    inherit the observation that an intention to communicate is itself a
    representational state, so that thought is presupposed prior to language.%
    \unverified{The thought-priority objection to intention-based accounts needs a
    citation. A candidate is Devitt, ``Three Mistakes About Semantic Intentions''.}
    Tomasello's account is explicit on this point. It posits shared intentionality,
    joint attention and cooperative motives as a biological infrastructure that evolved
    before conventional language and made it possible \cite{Tomasello2008}. The question
    of reference is moved to the evolution of that infrastructure, not answered from
    nothing.
  \item \emph{Computational accounts.} Berwick and Chomsky locate the decisive innovation
    in a single combinatorial operation, Merge, and explain how hierarchical structure
    is built from lexical atoms \cite{BerwickChomsky2016}. They accept that words
    preceded grammar, and critics have pointed out that the origin of those atoms is left
    largely unaddressed \cite{StuddertKennedyTerrace2017}.
\end{itemize}

None of these observations is a refutation. Each account addresses its own question, and
every model must take something as given, including the one developed here, which
assumes perception and response. The point concerns where the line is drawn. In each
case, the step by which some piece of behaviour first comes to \emph{stand for} something
is not part of what is explained.

\paragraph{Evidence.}
Language leaves no fossils. The events in question are prehistoric, unobservable and
unrepeatable, and the indirect evidence from anatomy, archaeology, comparative cognition
and genetics constrains them only weakly. In 1866 the Soci\'et\'e de Linguistique de
Paris went as far as to exclude communications on the origin of language from its
proceedings \cite{Lyons1991}. Much more recently, a review by leading researchers
concluded that the central questions remain largely out of reach of the available
evidence \cite{Hauser2014}. Theories of origin are therefore often regarded as
untestable, and dismissed as just-so stories: narratives that fit what they explain, but
that no observation could refute.

These two problems point toward a particular kind of contribution. A theory of origin
that cannot be tested historically may still be assessed in another way: by asking
whether its conclusions \emph{follow} from its assumptions, and which assumptions they
require. A claim of the form ``if these conditions hold, this structure can arise'', or
``must have this form'', is not a claim about what happened. It can be checked
deductively, and its assumptions can then be examined individually, some of them
empirically, for plausibility. This does not make an origin story testable as history. It
makes it accountable as an argument.

\paragraph{The originary hypothesis.}
Eric Gans's Generative Anthropology is notable for addressing the first problem
directly, and in a way that lends itself to this kind of assessment
\cite{Gans1981,Gans1993}. Several features distinguish it.
\begin{enumerate}
  \item It starts from the sign itself. It asks what must happen for a sign to exist at
    all, not how pre-existing contents come to be communicated.
  \item It derives the sign from non-semiotic behaviour. On Gans's account, the first
    sign is an aborted gesture of appropriation: a practical act that, by being
    interrupted, comes to designate instead of to seize. This gives a concrete mechanism
    by which representation could arise without being presupposed.
  \item The sign is shared from the outset. It arises in a scene among several
    participants, as a mutually understood gesture, so meaning is not a private state
    later made public.
  \item It is generative. A single minimal hypothesis is followed by a derivation of
    further elementary forms of language in sequence. A claimed derivation has internal
    structure that can be examined step by step.
\end{enumerate}
\emph{The Origin of Language} opens with the sentence ``Mysteries should not be
multiplied beyond necessity'' \cite{Gans1981}. It is a version of Ockham's razor, and it
states the ambition of the theory: to explain the origin of the human while assuming as
little as possible.

The hypothesis runs as follows. A group of hominids converge, through mimetic desire, on
a single desirable object. Mutual imitation intensifies their rivalry beyond what an
animal dominance hierarchy can contain,%
\unverified{The role of the dominance hierarchy in Gans's scene has not been confirmed
in a source. Check Gans 1981 or 1993.}
and the group faces a violent crisis. The crisis is deferred when the gesture of
appropriation is aborted. The aborted gesture, shared and reciprocally understood,
becomes the first sign: an \emph{ostensive} that designates the object as a centre which
no one may take. The scene concludes in the \emph{sparagmos}, the collective division of
the object.

From the ostensive, Gans derives the further forms in sequence. The \emph{imperative}
arises from an ostensive made in the absence of its referent, which the hearer answers by
supplying the missing object. The \emph{declarative} arises when the imperative cannot be
fulfilled, and the hearer responds with a negation of the object's presence.%
\unverified{This summary of the imperative (the ``inappropriate ostensive'') and the
declarative (the negative ostensive) was confirmed only through secondary commentary.
Check against Gans 1981.}

Later work in the tradition develops related notions. Adam Katz, who also writes as
Dennis Bouvard, introduced into the originary scene the \emph{firstness} of a single
participant who emits the sign before the others \cite{Katz2005}; see also
\cite{Gans2007}. He has also developed an approach organized around the \emph{centre}
toward which the attention of a scene converges.%
\unverified{Katz on the centre has not been read directly. A candidate source is Katz,
``Generative Anthropology as the One Big Discipline'', \emph{Anthropoetics} XXIV, no.~2
(2019).}

\subsection{The problem: narrative content of uncertain necessity}

The present work grows out of a long engagement with the writings of Gans and Katz, and
out of a specific dissatisfaction with one feature of the account. Gans's hypothesis
avoids presupposing the sign. The scene that does this work, however, carries
substantial narrative content:
\begin{itemize}[nosep]
  \item \emph{convergence}: mimetic desire draws the participants to a single object;
  \item \emph{crisis}: mimetic rivalry exceeds what an animal dominance hierarchy can
    contain;
  \item \emph{deferral}: the prospect of violence motivates the abortion of the gesture
    of appropriation;
  \item \emph{centre}: the designated object becomes a sacred centre;
  \item \emph{resolution}: the scene ends in the sparagmos, in which the tension the sign
    has deferred is discharged on the object, not on other members of the group.%
    \unverified{This description of the sparagmos comes from an \emph{Anthropoetics}
    page whose author has not been confirmed. Find a citable source.}
\end{itemize}
Much of this content is inherited from Ren\'e Girard's mimetic theory
\cite{Girard1977}. Gans accepts that language originates in a mimetic crisis, while
rejecting the scapegoat mechanism as sufficient to produce a sign. The sparagmos is in
part his way of incorporating Girard's account of sacrificial violence.%
\unverified{The relation between Gans and Girard summarized here has been confirmed
only through an encyclopedia entry. Find a primary or scholarly source.}
Each element is anthropologically suggestive, and each ties the origin of language to
themes the theory develops elsewhere, such as the sacred, ritual, ethics and exchange.

Which of these elements is actually needed is a natural question for an account that
begins by refusing to multiply mysteries. It has been raised within the tradition itself.
Sackeroff, for example, identifies the mechanism of deferral as a point of difficulty,
observing that the account draws on two different sources of motivation: animal
appetite for the object, and a distinctively human investment in the sign
\cite{Sackeroff2015}.%
\unverified{The issue number of Sackeroff's article was inferred from its URL. Confirm.}
The present work takes up this internal question in a formal setting.

Several of the scene's elements are open to objection on their own terms.
\begin{itemize}
  \item \emph{Violence.} Among primates, conflict is frequent but usually limited, and it
    is often followed by reconciliation \cite{deWaal1989}. Evolutionary game theory
    explains why contests between members of the same species usually take the form of
    limited, not escalated, fighting \cite{MaynardSmithPrice1973}. Disputes that end in
    exhaustion, hierarchy or reconciliation do not by themselves produce signs. The
    account's answer is that mimetic intensity has disabled these ordinary resolutions,
    which is itself one of the elements whose necessity is in question.
  \item \emph{Deferral as a strategy.} In the Hawk--Dove game, mutual restraint is not
    stable: against a participant who refrains, taking the resource does better.%
    \unverified{The Hawk--Dove game postdates Maynard Smith and Price (1973). Cite the
    right source, probably Maynard Smith, \emph{Evolution and the Theory of Games}
    (1982).}
    A deferral that rests on self-interest therefore needs further explanation. The
    account supplies it through the intensity of mimetic desire and the threat of
    collective violence, again through the narrative content.
  \item \emph{Division.} By the time the sparagmos begins, the aborted gesture has
    already designated the object. Its role in the origin of distribution and of the
    sacred is not in question here. Its role in the emergence of reference is unclear.
\end{itemize}

The issue is therefore not whether these elements are plausible. It is whether they are
\emph{load-bearing}. If the derivation of the sign and its forms depends on them, it
inherits their contingency, and the charge of a just-so story applies. If it does not,
the derivation should be statable without them.

The present work follows in the spirit of Gans's opening sentence, and applies it to the
originary scene itself. Mimetic rivalry, catastrophic violence and the sparagmos are
explanatory commitments, and each is a candidate for a mystery multiplied beyond
necessity. The question is whether the derivation of the sign and its forms requires
them. If it does not, the principle with which Gans begins recommends setting them aside.

\subsection{The target: a grounded symbolic graph}

A generative account needs a destination as well as a starting point. Both can be stated
using accepted notions of the sign.

\paragraph{The starting state.}
There are actors and physical objects. There are contexts: at each time, which actors are
present and which objects are within whose reach. Actors move, and some movements are
directed at objects. All the relations in this state are momentary and local. Each holds
between particular actors and things, at a particular time, in a particular context.
There are no symbolic nodes.

\paragraph{The sign as a relation.}
The accepted definitions of the sign are relational. For Saussure, the linguistic sign
joins a concept and an acoustic image, and its value is fixed only by its differences
from other signs in the system.%
\unverified{Add a citation, with page, to the \emph{Course in General Linguistics}.}
For Peirce, ``a sign, or representamen, is something which stands to somebody for
something in some respect or capacity'' \cite[2.228]{PeirceCP}. This is a relation among
a sign, an object and an interpretant, and it holds \emph{for} someone. Morris divides
the study of signs by the relations a sign can enter into: to other signs, to objects,
and to its interpreters \cite{Morris1938}.

On each account, a sign is not a thing with a property but a term in relations. Among
those relations are relations to the \emph{actors} who produce and interpret it. The
minimal mathematical structure of relations is a graph. Relations among more than two
terms can be represented as a hypergraph, or encoded in an ordinary graph by an incidence
construction, in which each instance of a relation becomes a node joined to its
arguments. The relevant graphs are many-sorted: their vertices include actors, physical
objects, and the acts by which actors relate to objects and to one another.

\paragraph{Actors, time and context.}
Every such relation is indexed by actors, by a time and by a context. The most
elementary sign relation, ostension, shows all three at once. \emph{This}, \emph{here} and
\emph{now} designate only within the situation in which they are produced, for those who
share it. A pointing gesture relates one actor, by way of a movement, to an object, at one
moment and in one place, and only for another actor who takes it up. The relation does
not outlast the act.

Symbolic reference is what escapes this indexing. Hockett lists \emph{displacement}, the
capacity to refer to what is removed in time or place from the moment of speaking, among
the design features of human language \cite{Hockett1960}. Deacon, building on Peirce,
argues that symbolic reference is grounded in indexical reference but depends on
relations among symbols themselves \cite{Deacon1997}. In graph terms, a symbolic node is
something that persists across contexts and across occasions of use: different actors
can re-identify it, at other times and in other places, and it can enter into further
relations there.

\paragraph{The end state.}
A generative account must arrive at a graph with the following properties.
\begin{enumerate}
  \item \emph{Nodes that stand for things.} Some vertices are symbolic, and stand for
    objects to actors.
  \item \emph{Persistence beyond the context of origin.} A symbolic node can be
    re-identified at another time, in another place and by another actor, including in
    the absence of what it stands for.
  \item \emph{Relations among symbolic nodes.} Symbolic nodes stand in relations to one
    another, not only to objects, and some of these relations have \emph{roles}, where it
    matters which term is which.
  \item \emph{Grounding.} Every symbolic node can be traced back to the physical world,
    so there are no free-floating symbols.
  \item \emph{Truth and falsity.} What is said with symbolic nodes has conditions under
    which it holds, and it can fail.%
    \unverified{Hockett and Altmann's ``prevarication'' may be the right citation here.
    Check.}
  \item \emph{Partiality.} Relations can be affirmed, denied, or left unsaid, without
    silence having to carry information.
  \item \emph{Open-endedness.} Any finite relational structure over grounded objects can
    be expressed.
\end{enumerate}

\paragraph{The generative task.}
The task is to get from the starting state to the end state using only what the starting
state contains, and operations on it applied to their own products. Nothing symbolic may
be assumed along the way. Several features commonly attributed to human language lie
outside the target: the vocal-auditory channel, duality of patterning, and cultural
transmission across generations. The account addresses the emergence of reference and of
relational structure, not of speech.

The target also makes precise a hypothesis that motivated the work:
\begin{itemize}[nosep]
  \item the ostensive establishes an edge between actors and an object within a single
    context;
  \item the imperative is such an edge whose object is missing from the actors' current
    context;
  \item the declarative yields a node: a relation that has itself become something the
    actors can refer to, and that persists across contexts.
\end{itemize}
The formalization supports a sharper version of the last point. The first symbolic node
is a \emph{reified relation}, and relations between nodes arise by the same operation
applied one level up.

\subsection{Aim: a minimal axiom set that does not license itself}

The aim is to find a minimal set of axioms under which the end state of
Section~1.3 can be reached from the starting state, and to prove that it can.

\paragraph{The rule against self-licensing.}
An axiom set \emph{licenses itself} when what it is supposed to generate is built into its
assumptions, so that the derivation only returns what was put in. To exclude this, every
axiom set considered here must satisfy two conditions.
\begin{enumerate}
  \item \emph{Base vocabulary only.} Axioms may speak only of what the starting state
    contains: actors, physical objects, times, movements, reach, contact, perception, and
    resemblance between movements. They may not mention signs, meaning, reference,
    uptake, the ostensive, the imperative, the declarative, symbolic nodes, claims or
    truth, nor any notion defined from these.
  \item \emph{Target notions by definition only.} Every notion that belongs to the end
    state is introduced by an explicit definition in the base vocabulary. A definition
    adds no constraint on what can happen: it names a pattern, which may or may not
    occur.
\end{enumerate}
Under these conditions, every property of the end state that the account establishes is a
theorem, derived from axioms that do not mention it.

\paragraph{Definitions are hypotheses.}
The rule shifts the weight of the account onto its definitions: what it is for a movement
to be taken up as being about something, or for a pattern of responses to constitute a
shared centre. A definition can smuggle in as much as an axiom can. Two safeguards apply.
Definitions are drawn as closely as possible from accepted notions of the sign, and they
are stated as explicit hypotheses about what the defined notion consists in, open to
dispute and to replacement by alternatives.

\paragraph{Generativity relative to the target.}
With the target fixed and the rule in force, the aim has three parts.
\begin{description}[leftmargin=1.5em, style=nextline]
  \item[Existence.] There is a model of the axioms in which the end state, or the relevant
    part of it, arises from a starting state that lacks it.
  \item[Necessity.] In every model of the axioms, whenever part of the end state exists,
    it has a determinate origin and order.
  \item[Universality.] The operations that produce the first symbolic structures suffice
    for open-endedness: any finite relational structure over grounded objects can be
    reached.
\end{description}

\paragraph{Minimality.}
An axiom set is minimal when each axiom is needed: for each axiom there is a model of the
others in which some property of the end state fails. Minimality is shown axiom by axiom,
by exhibiting such countermodels.

The aim is not to show that the end state was inevitable, nor that it was reached in one
particular way. It is to show that it can be reached under assumptions that are weak,
explicit and free of the target, and to identify which assumptions each part of it
requires.

\paragraph{Comparing accounts.}
The same rule suggests a way of comparing accounts. Formalize an account, separate its
primitives into base vocabulary and anything semiotic, and ask whether its target notions
can be defined from the base part alone. Padoa's method from logic gives a standard test:
a notion is not definable from the other primitives if two models agree on all of them
but disagree on that notion. One feature that may signal self-licensing is a primitive
kind of act whose only role is to be responded to, such as the costless signals of a
signaling game. Any such result holds only relative to a particular formalization, and
only relative to the present target, the emergence of sign-hood. It is not a criticism
of what other accounts set out to explain. Applying the test to other accounts,
including Gans's own scene, is left as work in progress.

\subsection{Rival orderings}

A further motivation is comparative. The sequence ostensive, imperative, declarative is
one reading among several, and other readings, for instance ones in which the
declarative is primary, are possible.%
\unverified{Name proponents of declarative-first or other orderings, or rephrase as the
author's conjecture.}
Under the rule of Section~1.4, the choice between them becomes precise. An account that
builds a particular mode into its axioms licenses its own ordering. The real question is
which orderings can be generated from axioms in the base vocabulary, and under which
assumptions. The framework developed here bears on this directly. In it, the declarative
requires the \emph{absence} of its object, but not a prior failed imperative. The
imperative and the declarative appear as two responses to absence, not as successive
stages.

\subsection{Why machine verification}

Arguments in humanistic theory often rely on steps that are intuitively compelling but
unstated. Formalization in a proof assistant forces every assumption to be written down
and every inference to be checked. All results reported here are formalized in the
Lean~4 proof assistant \cite{deMoura2021} and checked in full, with no unproven steps.

The process was itself informative. Several assumptions that the informal narrative
leaves implicit were exposed only when a proof failed without them. Among them:
\begin{itemize}[nosep]
  \item a kinematic principle linking a reproduced gesture to the object it was
    originally directed at;
  \item the requirement that one participant act before another, which corresponds to
    Katz's firstness;
  \item the dependence of falsehood on the drift of signs away from their iconic
    origins;
  \item the reliance of affirmation, derived as the negation of a negation, on classical
    logic.
\end{itemize}
Formalization also makes it possible to say exactly which statements are derived and
which are merely definitions. The paper marks this distinction throughout.

\subsection{Outline}

Section~2 introduces the primitives and axioms, which concern only actors, movements,
reach and resemblance. Section~3 defines taking-as, the centre, dependent response and the
elementary modes as patterns of response. Section~4 reports the results that hold in every
model: the origin of signs, the form of the imperative, the declarative's dependence on
absence, direction of fit, and the conditions for falsehood. Section~5 reports the results
on graphs: claims, universality, roles and polarity. Section~6 reports the models and
countermodels that establish existence and show which assumptions are needed. Section~7
discusses the relation to Generative Anthropology, open problems, and directions for
collaborative work. Table~\ref{tab:roadmap} shows where each property of the end state is
reached.

\begin{table}[h]
\centering
\small
\begin{tabular}{@{}p{0.36\linewidth}p{0.6\linewidth}@{}}
\toprule
Property of the end state & Where it is reached \\
\midrule
1. Nodes that stand for things & declarative and lineage (Section 3) \\
2. Persistence & re-presentation; displaced claims (Sections 4--5) \\
3. Relations, with roles & claims and ordered claims (Section 5) \\
4. Grounding & content traces to an object once within reach (Section 4) \\
5. Truth and falsity & falsehood requires drift (Section 4) \\
6. Partiality & polarity; partial universality (Section 5) \\
7. Open-endedness & universality (Section 5) \\
\bottomrule
\end{tabular}
\caption{Where each property of the end state is reached.}
\label{tab:roadmap}
\end{table}

\begin{thebibliography}{99}

\bibitem{AlexanderSkyrmsZabell2012}
J.~M. Alexander, B.~Skyrms and S.~L. Zabell.
\newblock Inventing new signals.
\newblock \emph{Dynamic Games and Applications}, 2012.
\newblock doi:10.1007/s13235-011-0027-2.

\bibitem{Arbib2012}
M.~A. Arbib.
\newblock \emph{How the Brain Got Language: The Mirror System Hypothesis}.
\newblock Oxford University Press, New York, 2012.

\bibitem{BerwickChomsky2016}
R.~C. Berwick and N.~Chomsky.
\newblock \emph{Why Only Us: Language and Evolution}.
\newblock MIT Press, Cambridge, MA, 2016.

\bibitem{Corballis2002}
M.~C. Corballis.
\newblock \emph{From Hand to Mouth: The Origins of Language}.
\newblock Princeton University Press, Princeton, 2002.

\bibitem{deMoura2021}
L.~de Moura and S.~Ullrich.
\newblock The Lean 4 theorem prover and programming language.
\newblock In \emph{Automated Deduction -- CADE 28}, Lecture Notes in Computer
  Science 12699, pages 625--635. Springer, 2021.

\bibitem{deWaal1989}
F.~de Waal.
\newblock \emph{Peacemaking among Primates}.
\newblock Harvard University Press, Cambridge, MA, 1989.

\bibitem{Deacon1997}
T.~W. Deacon.
\newblock \emph{The Symbolic Species: The Co-evolution of Language and the Brain}.
\newblock W.~W. Norton, New York, 1997.

\bibitem{Dunbar1996}
R.~Dunbar.
\newblock \emph{Grooming, Gossip, and the Evolution of Language}.
\newblock Harvard University Press, Cambridge, MA, 1996.

\bibitem{Gans1981}
E.~Gans.
\newblock \emph{The Origin of Language: A Formal Theory of Representation}.
\newblock University of California Press, Berkeley, 1981.
\newblock New edition: \emph{The Origin of Language: A New Edition}, with an
  introduction by A.~Katz, Spuyten Duyvil, New York, 2019.

\bibitem{Gans1993}
E.~Gans.
\newblock \emph{Originary Thinking: Elements of Generative Anthropology}.
\newblock Stanford University Press, Stanford, 1993.

\bibitem{Gans2007}
E.~Gans.
\newblock On firstness.
\newblock In A.~Katz, editor, \emph{The Originary Hypothesis: A Minimal Proposal for
  Humanistic Inquiry}, pages 45--57. The Davies Group, Aurora, CO, 2007.

\bibitem{Girard1977}
R.~Girard.
\newblock \emph{Violence and the Sacred}.
\newblock Translated by P.~Gregory. Johns Hopkins University Press, Baltimore, 1977.

\bibitem{Grice1957}
H.~P. Grice.
\newblock Meaning.
\newblock \emph{The Philosophical Review}, 66(3):377--388, 1957.

\bibitem{Hauser2014}
M.~D. Hauser, C.~Yang, R.~C. Berwick, I.~Tattersall, M.~J. Ryan, J.~Watumull,
  N.~Chomsky and R.~C. Lewontin.
\newblock The mystery of language evolution.
\newblock \emph{Frontiers in Psychology}, 5:401, 2014.

\bibitem{Hockett1960}
C.~F. Hockett.
\newblock The origin of speech.
\newblock \emph{Scientific American}, September 1960.

\bibitem{Katz2005}
A.~Katz.
\newblock Remembering Amalek: 9/11 and generative thinking.
\newblock \emph{Anthropoetics}, X(2), Fall 2004/Winter 2005.

\bibitem{Lewis1969}
D.~Lewis.
\newblock \emph{Convention: A Philosophical Study}.
\newblock Harvard University Press, Cambridge, MA, 1969.

\bibitem{Lyons1991}
J.~Lyons.
\newblock The origin of language, speech and languages.
\newblock In \emph{Natural Language and Universal Grammar: Essays in Linguistic
  Theory}, volume~1. Cambridge University Press, Cambridge, 1991.

\bibitem{MaynardSmithPrice1973}
J.~Maynard Smith and G.~R. Price.
\newblock The logic of animal conflict.
\newblock \emph{Nature}, 246:15--18, 1973.

\bibitem{Morris1938}
C.~W. Morris.
\newblock \emph{Foundations of the Theory of Signs}.
\newblock International Encyclopedia of Unified Science. University of Chicago Press,
  Chicago, 1938.

\bibitem{PeirceCP}
C.~S. Peirce.
\newblock \emph{Collected Papers of Charles Sanders Peirce}.
\newblock Harvard University Press, Cambridge, MA. Cited by volume and paragraph.

\bibitem{Sackeroff2015}
S.~Sackeroff.
\newblock The ends of deferral.
\newblock \emph{Anthropoetics}, XX(2), 2015.

\bibitem{Skyrms2010}
B.~Skyrms.
\newblock \emph{Signals: Evolution, Learning, and Information}.
\newblock Oxford University Press, Oxford, 2010.

\bibitem{StuddertKennedyTerrace2017}
M.~Studdert-Kennedy and H.~Terrace.
\newblock In the beginning: a review of Robert C. Berwick and Noam Chomsky's
  \emph{Why Only Us}.
\newblock \emph{Journal of Language Evolution}, 2(2):114--125, 2017.

\bibitem{Tomasello2008}
M.~Tomasello.
\newblock \emph{Origins of Human Communication}.
\newblock MIT Press, Cambridge, MA, 2008.

\end{thebibliography}

\end{document}
